\section{Related Work}
\label{sec:related_work}

Building a plant disease classifier that is accurate, edge-deployable, explainable, and reliable under real field conditions draws on several largely separate literatures. We review each in turn: the datasets and CNN-based classifiers that establish the task (Section~\ref{sec:related_work}.1), lightweight architectures for edge deployment (Section~\ref{sec:related_work}.2), pixel-level and language-based explanation methods (Sections~\ref{sec:related_work}.3--\ref{sec:related_work}.4), and the data-efficiency and reliability concerns that determine whether such a classifier can be trusted in practice (Sections~\ref{sec:related_work}.5--\ref{sec:related_work}.7).

\subsection{Plant Disease Classification}
Early works on plant disease recognition used traditional machine learning with hand‑crafted features; the introduction of deep learning, especially convolutional neural networks (CNNs), has since become the dominant paradigm. The PlantVillage dataset \cite{hughes2015plant} contains more than 50,000 images of healthy and diseased leaves under controlled conditions and underlies a large share of this literature, with models such as VGG, ResNet, and Inception routinely reported near ceiling accuracy on it \cite{mohanty2016using}. Two recent surveys covering work from 2015 through 2023 confirm this pattern at scale, and both single out limited robustness under field conditions as the field's central open problem, since PlantVillage-level accuracy is understood to reflect the dataset's controlled imaging rather than a solved recognition problem \cite{shoaib2023review, mustofa2023review}. PlantDoc \cite{singh2020plantdoc}, a smaller (2,598-image), internet-sourced, field-style dataset spanning 13 species on which PlantVillage-pretrained models transfer poorly, gives this concern a concrete empirical footing. Field‑collected datasets nonetheless remain scarce more broadly \cite{barbedo2019plant}, and avocado-specific datasets are especially rare. The only one we are aware of, K-Kotagiri \cite{kotagiri2024avocado}, is a 435-image, binary (healthy versus diseased) dataset collected in Tamil Nadu, India, without disease-level labels. AvocadoLeaf is positioned squarely in this gap: 1,861 images spanning five disease categories and a healthy class, captured under diverse natural conditions, including variable lighting, multiple angles, and cluttered backgrounds. Table~\ref{tab:dataset_comparison} places it alongside PlantVillage, PlantDoc, and K-Kotagiri along the dimensions that matter here, imaging setting and species/class coverage, rather than raw image count; K-Kotagiri's independent collection site and labelling scheme also make it the candidate dataset for the cross-domain evaluation in Section~\ref{sec:experiments}.

\begin{table*}[htbp]
\centering
\caption{Comparison of AvocadoLeaf with three related plant disease benchmarks.}
\label{tab:dataset_comparison}
\small
\begin{tabularx}{\textwidth}{l c c X}
\toprule
\textbf{Dataset} & \textbf{Images} & \textbf{Species / classes} & \textbf{Setting} \\
\midrule
PlantVillage \cite{hughes2015plant} & $>$50,000 & multi-crop & lab-controlled (uniform background, fixed illumination) \\
PlantDoc \cite{singh2020plantdoc} & 2,598 & 13 species & field, internet-sourced (uncontrolled and uncurated) \\
K-Kotagiri \cite{kotagiri2024avocado} & 435 & 1 species (avocado), 2 classes (healthy/diseased) & field-collected, single site (Tamil Nadu, India), no disease-level labels \\
AvocadoLeaf (ours) & 1,861 & 1 species (avocado), 6 classes & field-collected first-hand across 3 provinces (deliberate coverage of light, weather, angle) \\
\bottomrule
\end{tabularx}
\end{table*}

The trade-off is that PlantDoc's broader 13-species scope comes at the cost of far fewer images per class than a single-crop dataset such as AvocadoLeaf can provide (Section~\ref{sec:dataset_construction}). K-Kotagiri, meanwhile, targets the same species but only at the coarse healthy/diseased level, which is why we treat it in Section~\ref{sec:experiments} as a cross-domain \emph{binary} transfer test rather than a like-for-like replacement for AvocadoLeaf's 6-class taxonomy.

\subsection{Lightweight Architectures for Edge Deployment}
Deploying deep learning models on resource‑constrained edge devices (e.g., smartphones, Raspberry Pi) requires careful trade‑offs between accuracy, model size, and inference speed. MobileNet \cite{howard2017mobilenets} and EfficientNet \cite{tan2019efficientnet} are popular families of lightweight CNNs that use depthwise separable convolutions and neural architecture search, while ConvNeXt \cite{liu2022convnet}, and its V2 revision \cite{woo2023convnextv2} used as the proposed model in this paper (Section~\ref{sec:proposed_model}), instead modernise standard CNNs with design elements borrowed from vision transformers to reach a favourable accuracy‑efficiency trade‑off. Applications of lightweight models to plant disease classification typically adopt one such family and benchmark a handful of its variants \cite{atila2021plant}, rather than comparing across CNN, hybrid CNN‑Transformer, and pure transformer families under one shared protocol. Section~\ref{sec:experiments} takes the broader view, benchmarking 15 architectures spanning all three families to identify ConvNeXt‑V2 as offering the best accuracy‑efficiency trade‑off for edge deployment. The shared protocol has one deliberate exception (the ResNet-18 baseline unfreezes one stage rather than two, Section~\ref{sec:proposed_model}) and one undisclosed one (an augmentation-policy inconsistency surfaced only after training, Section~\ref{subsec:ablation}); Section~\ref{sec:discussion} discusses both and why neither is judged likely to change which architecture is selected; Sections~\ref{subsec:cross_dataset} and~\ref{subsec:realtime_demo} then go beyond the batched-GPU-latency comparisons typical of this literature, probing the chosen architecture's behaviour under distribution shift and its live inference latency on a laptop webcam, respectively.

\subsection{Explainable AI for Plant Diagnosis}
Understanding why a model makes a particular prediction is crucial for building trust with end‑users. Local Interpretable Model-agnostic Explanations (LIME) \cite{ribeiro2016why} is a popular model-agnostic technique for visualising the superpixels of an input image that most influence the model's decision, while gradient-based methods such as Grad-CAM \cite{selvaraju2017grad} instead produce a coarser class-activation heatmap directly from a CNN's internal gradients. In agriculture, XAI has been used to highlight disease symptoms on leaves \cite{sharma2020explainable}, helping farmers verify model outputs, though typically through qualitative visualisation on a handful of hand-picked examples rather than a quantitative measure of how faithful the highlighted region is to the model's actual decision. Section~\ref{sec:experiments} closes this gap for LIME specifically, applying a deletion metric across the full 280-image test set to test, image by image, whether removing the highlighted region collapses the model's confidence as a faithful explanation should.

\subsection{Vision-Language Models for Agricultural Explanation}
Beyond pixel-level attribution, a fast-growing body of work explores large vision-language models (VLMs) for plant diagnosis. One line of work fine-tunes a VLM as the diagnostic decision-maker itself. AgroGPT \cite{awais2025agrogpt} builds an instruction-tuning dataset from agricultural vision-only data to support conversational, fine-grained crop and disease identification, and AgriChain \cite{mahmood2026agrichain} fine-tunes a VLM on expert-verified, image-grounded diagnostic rationales so the same model predicts the class and produces its supporting explanation. A second, benchmark-oriented line evaluates VLMs on large visual-question-answering (VQA) corpora built specifically for plant science, such as PlantExpertVQA \cite{sakib2025plantexpertvqa} (over 765,000 question-answer pairs across 38 crop species) and a two-stage crop-disease VQA framework that pairs a frozen vision encoder with Grad-CAM and token-level attribution for interpretability \cite{hossain2026cropvqa}. Across both lines, the VLM's own classification or generation ability is what ultimately drives correctness. In contrast, our pipeline keeps a separately validated CNN as the sole decision-maker and uses a small, free, locally-hosted VLM (\texttt{Qwen2-VL-2B-Instruct}, a similar parameter scale to the compact models fine-tuned in \cite{sakib2025plantexpertvqa}) only to narrate the CNN's already-decided prediction together with LIME's visual attribution (Section~\ref{subsec:vlm_explain}). We additionally report, as a negative result, that letting the same VLM's judgement double as an automatic error-detection signal for the CNN does not outperform the CNN's own predictive-entropy baseline (\ref{app:llm}).

\subsection{Few-shot Learning for Plant Disease Recognition}
Collecting large annotated datasets for every possible disease is impractical, which motivates few‑shot learning: recognising new classes from only a handful of labelled examples. Popular approaches include metric‑based methods such as Prototypical Networks (ProtoNet) \cite{snell2017prototypical} and Relation Networks \cite{sung2018learning}, as well as optimisation‑based methods like MAML \cite{finn2017model}, and in plant pathology such methods have been applied to tomato, rice, and cassava diseases \cite{arguelles2020few, wang2021few}. These applications typically still fine-tune or meta-train an ImageNet-pretrained backbone on the target disease domain, rather than testing a frozen, already-trained classification backbone directly as a metric-learning feature extractor. Section~\ref{subsec:fewshot} takes the latter route, evaluating ProtoNet and Relation Networks with a frozen ConvNeXt‑V2 backbone (already trained for full-shot classification, Section~\ref{sec:experiments}) on the fine‑grained AvocadoLeaf classes, and compares both against conventional fine‑tuning under the same shot budgets.

\subsection{Uncertainty Estimation in Deep Learning}
High accuracy alone is insufficient for high‑stakes applications; a reliable model must also be well calibrated and able to detect out‑of‑distribution (OOD) inputs. Expected Calibration Error (ECE) \cite{guo2017calibration} is a standard metric for assessing calibration. Monte Carlo (MC) Dropout \cite{gal2016dropout} provides a simple yet effective way to estimate predictive uncertainty by treating dropout at inference time as approximate Bayesian sampling, requiring only one trained model. Deep Ensembles \cite{lakshminarayanan2017ensembles} instead train several independent models from different random initialisations and combine their predictions, typically yielding better-calibrated uncertainty at the cost of a multiplicative increase in training and storage cost. We adopt MC Dropout over ensembling for this reason, since it reuses the single ConvNeXt‑V2 model already trained for classification (Section~\ref{sec:experiments}) rather than requiring several independently trained copies. Calibration \cite{wang2022uncertainty} and OOD detection \cite{gonzalez2020identifying} have each received some attention in plant disease classification, but rarely together, and rarely for a model already selected for edge deployment on efficiency grounds. Section~\ref{sec:experiments} reports both for the same lightweight avocado-disease model, measuring calibration directly and testing OOD detection against two distinct kinds of shift, far (random internet images) and near (leaves of other crop species), so that the two aspects of reliability can be read against the same accuracy and efficiency numbers established earlier in that section.

\subsection{Cross-dataset Generalization in Plant Disease Classification}
A model's accuracy on its own held-out test set does not guarantee it will generalise to images from a different source. Wu et al.\ \cite{wu2023laboratory} show that networks trained on laboratory-style images suffer marked accuracy loss when applied to field-collected images of the same diseases, and propose unsupervised domain adaptation to narrow this laboratory-to-field gap. Sun \cite{sun2024crossdataset} instead trains standard architectures (ResNet, EfficientNet, Inception, MobileNet) on one tomato pest/disease dataset and evaluates them directly on independently collected tomato datasets with no adaptation step, finding substantial accuracy degradation regardless of architecture; Yang et al.\ \cite{yang2024laboratory} report a similar lab-to-field gap and propose cross-domain few-shot learning to narrow it, evaluating on their own independently collected Field Crop Disease Dataset. All three studies rely on a second, independently labelled dataset covering the same classes to measure this gap directly as a classification-accuracy drop. We do not yet have such a second AvocadoLeaf-like dataset (Section~\ref{sec:discussion}); instead, Section~\ref{subsec:cross_dataset} reuses the far- and near-OOD image sets already collected for out-of-distribution detection to characterise a related but distinct question, namely how confidently a naive (non-uncertainty-gated) deployment would misclassify domain-shifted, non-avocado inputs. Collecting an independent, same-task dataset to enable a true cross-dataset accuracy evaluation in the sense of \cite{wu2023laboratory, sun2024crossdataset, yang2024laboratory} is left as future work.
